\subsection{Notation for the limit and the board}\label{sec:notacion}
Table~\ref{tab:notacion} distinguishes the objects on the finite board from the quantities associated
with the profile. In particular, a marked pair and the matrix whose minor is taken appear together in
Appendix~\ref{app:boundary}.
\begin{center}
\small
\begin{tabular}{@{}L{.21\linewidth}L{.72\linewidth}@{}}
\toprule Symbol & Meaning\\ \midrule
\(\Phi,\ C_{\rm pred}\) & Coefficient of \(n\) and constant term of the original predictor.\\
\(R_Q,\ \cerr\) & Difference of the two logarithms and limit of \(nR_Q(n)\).\\
\(m,\ j,\ x\) & Empty rows at a step, arrival size \(\lfloor bn\rfloor\), and residual fraction \(m/n\).
In the rook lemmas \(k\) is the size of a mask.\\
\(B_m,\ M_m,\ P\) & Residual board, rook mask, and doubly stochastic scaling of that mask.\\
\(\parmar,\ \familias(z)\) & Marked pair of squares and set of the four lines through \(z\).\\
\(H_j,\ F_b\) & Short historical arrival and arrival without an exit and with at least one completion.\\
\(J,\ S,\ \sumafrontera\) & Contribution of a single diagonal, magnitude of the square term, and signed
boundary sum \(\sumafrontera=P_T+J-S\); see~\eqref{eq:J-def} and~\eqref{eq:S-def}.\\
\(\coefT,\ F(x)\) & First rook coefficient, \(\coefT=\mathfrak F-1/12\), and value of \(\mathfrak F\) on
the scaled profile.\\
\(\centrado,\ \pesodirecto(z)\) & Centered rook kernel and weight of a crossing with its availability
indicators.\\
\(d(x),\ \numerr(x)\) & Sum of the sixteen contributions~\eqref{eq:d-suma} and numerator of the
integral in the theorem.\\
\(\valorcorte(b)\) & Approximation to the limit with the integration endpoint fixed at \(b\).\\
\bottomrule
\end{tabular}
\refstepcounter{table}\label{tab:notacion}\par\smallskip
\textbf{Table \thetable.} Symbols used in the proof.
\end{center}
